\chapter*{4. Datos}

Esta sección presenta el punto de partida para estimar estructuralmente el mercado de interés, centrándose en la construcción de una base de datos amplia y representativa de su historia y dinámica competitiva. El análisis riguroso de mercados de \textit{soft commodities} requiere capturar tanto las fluctuaciones de corto plazo como las tendencias estructurales, por lo que es fundamental integrar información histórica que incluya variables claves como precios, volúmenes de exportación y comercio, costos de producción, variables de control y factores climáticos. La selección de un horizonte temporal suficientemente extenso permite identificar patrones cíclicos, choques exógenos y relaciones de equilibrio, garantizando así que los resultados no solo reflejen elementos coyunturales, sino también la evolución estructural del sector.

En este caso, los datos recopilados resumen la información de 57 países exportadores y [] países importadores, para el período 1965-2019. [AGREGAR CAVEATS EN LA INFORMACION. TODOS LOS PAISES TIENEN MAXIMA COBERTURA?]


\subsection*{4.1. Datos utilizados para la estimación de demanda}

Como se explicará en la sección siguiente, la estimación de demanda se realiza a nivel agregado, por lo que se construyó una base de datos de frecuencia mensual con información de cantidades comercializadas y precios, complementados con variables de control e instrumentos que permitieron abordar el problema de endogeneidad [CITAR].

La información relativa a las cantidades exportadas fue obtenida del Departamento de Agricultura de los Estados Unidos (USDA), el cual reporta cifras anuales para cada país, incluyendo importaciones, exportaciones, nivel de producción (desagregado para Robusta y Arabica), consumo doméstico, \textit{stock} inicial y \textit{stock} final; todo expresado en millones de bolsas de 60 kg. Para esta investigación, se utilizó información de exportaciones de grano (ID 90) de cada país. Dado que la información solo está disponible en frecuencia anual, fue necesaria la aproximación mensual dividiendo las cantidades por 12. Este ajuste fue utilizado por Igami (CITAR) en la estimación estructural de la demanda, sin perjuicio de que también es posible llevar adelante un enfoque anual.

En cuanto a los precios, se utiliza un índice mensual proporcionado por la ICO, con base en información de precios \textit{ex dock}, que reflejan costos de producción, de transporte y por descarga de embarcaciones. Cuando esta información no está disponible, la ICO ajusta el índice FOB (\textit{Free on Borad}, el precio pre-desembarque), o bien el índice CIF (\textit{Cost, Insurance and Freight}, el precio post-desembarque), para emular el precio \textit{ex dock}. Luego, pondera los precios de cada variedad de café según su participación en las principales bolsas comerciales que transan este bien agrícola [CITAR]. Este índice, establecido en la década de 1960, se considera uno de los principales precios de referencia para la industria cafetalera mundial, utilizado para proyectar condiciones comerciales y evaluar futuros costos (\cite{ICO2025}).

Para controlar por el efecto ingreso de los países importadores, se incorpora el Producto Interno Bruto real. En particular, se adopta la medición por el lado del gasto del \textit{Penn World Table}, cuya metodología sigue a \cite{Feenstra2015}, y está expresada en dólares internacionales constantes de 2017, utilizando Paridades de Poder Adquisitivo (PPA) encadenadas, lo que permite comparaciones robustas del desempeño económico entre países y a lo largo del tiempo. Para la estimación de demanda, se utiliza la unidad expresada en logaritmos y se ajusta la frecuencia anual a una mensual, dividiendo por 12.

El efecto sustitución es controlado por el precio internacional del té. Se utiliza la serie obtenida de los \textit{World Bank Commodity Price Data} (conocidos como "\textit{The Pink Sheet}"). La serie consiste en precios mensuales en dólares estadounidenses nominales, disponibles desde 1960 hasta la actualidad. Cabe destacar que las series mensuales del Banco Mundial están disponibles únicamente en términos nominales y en dólares. En la estimación de demanda se utiliza el resultado ponderado del precio del té, el cual representa el promedio de las principales variedades de té en el mercado global.

Adicionalmente, dado que la estimación estructural de la demanda considera datos desde 1965, se utiliza una variable de control adicional para el período en que los países exportadores operaban bajo un acuerdo de cuotas.

Ahora bien, la determinación conjunta de precios y cantidades en el equilibrio de mercado presenta desafíos econométricos importantes, como son el sesgo por simultaneidad o la endogeneidad [CITAR]. Para abordar esta cuestión, se requiere la inclusión de variables instrumentales que influyan en la cantidad ofrecida sin afectar directamente la demanda del consumidor. En el contexto de los \textit{soft commodities}, las variables climáticas emergen como candidatos ideales, debido a su naturaleza exógena frente a las preferencias del consumidor y su impacto directo en la producción [CITAR]. La sección 4.3. describe con detalle la construcción de estas variables.

La información estadística de las variables requeridas para la estimación de demanda se resumen en la siguiente Tabla N°[].

[TABLA]


\subsection*{4.2. Datos utilizados para la estimación de costos marginales}

El objetivo principal de la estimación de costos marginales es determinar qué nivel de influencia tienen algunas variables sobre las condiciones de producción del café. En particular, interesa estimar el impacto de la exposición a temperaturas extremas, inadecuadas para el óptimo desarrollo de los cafetos, sobre los costos marginales implícitos que resultan de la estimación de demanda agregada. Para estos efectos, se utilizan diferentes \textit{cost-shifters}, con variación entre países y a nivel mensual.

En primer lugar, se utilizó un índice mensual de precios de fertilizantes, extraído del \textit{World Bank Commodity Price Data}, al igual que los datos del té. Este índice pondera series de precios de fertilizantes clave en la actividad agrícola, y se incluye como un determinante de los costos marginales de producción a nivel mensual, uniforme para todos los países.

Enseguida, se utiliza una variable que captura, para cada país y mes, la cantidad de días en que los campos de cultivo de café están expuestos a temperaturas sobre cierto umbral, definido a partir de la naturaleza del \textit{soft commodity} y las condiciones óptimas de temperatura para su correcto desarrollo. Por construcción, se trata de una variable continua, lo que permite preservar la variación completa en la intensidad del tratamiento climático y capturar efectos marginales en la relación entre estrés térmico y resultados económicos. El detalle de la construcción de esta variable se explica en la siguiente sección.

La Tabla N°[] resume la información de las variables revisadas.

[TABLA]

\subsection*{4.3. Variable de exposición a temperaturas extremas o estrés térmico}

La estimación estructural del mercado internacional de café requiere contar con una base de datos lo suficientemente extensa para abarcar el estrés térmico de cada país exportador y su evolución en el tiempo. Específicamente, la estimación de costos marginales demanda contar con datos de panel de frecuencia mensual, y la estimación de demanda requiere ponderar estos datos para llegar a un resultado agregado mundial.

Las variables de estrés térmico se construyen a partir de dos fuentes de datos clave. La primera de ellas otorga las temperaturas mínimas y máximas todo el mundo, lo que permite estimar la exposición diaria a temperaturas extremas. Por otro lado, la segunda fuente permite localizar espacialmente los cultivos de café dentro de cada país exportador, dando mayor precisión al análisis. En conjunto, estos datos aseguran que las mediciones reflejen las condiciones térmicas específicas que enfrenta el sector cafetalero en cada economía.

La base de datos del \textit{Inter-Sectoral Impact Model Intercomparison Project} (ISIMIP) proporciona información de temperaturas mínimas y máximas diarias con resolución espacial de 0.5° × 0.5° (55 km × 55 km en el ecuador, aproximadamente), ofreciendo una cobertura global consistente que es utilizada desde 1961 en adelante. Dentro de las bases disponibles, aquella con la designación \textit{obsclim} contiene datos observacionales corregidos y homogenizados, que combinan mediciones \textit{in situ} con reanálisis atmosférico, lo que garantiza la cobertura de variaciones climáticas localizadas para un bien cuyo desarrollo varían significativamente con las condiciones geográficas.

Por otro lado, la base de datos del \textit{Spatial Production Allocation Model} (SPAM 2020) proporciona información geoespacial detallada sobre la distribución global de áreas cosechadas para múltiples cultivos, incluyendo café Arabica y Robusta. Con una resolución espacial aproximada de 10 km × 10 km en el ecuador, SPAM desagrega la producción agrícola mediante un modelo de asignación que integra datos censales, estadísticas nacionales y variables auxiliares de clima, topografía y accesibilidad a mercados.

Las bases anteriores se complementan fuertemente, y es que la granularidad en la información de cada una permite identificar zonas específicas de cultivo de café dentro de cada país y revisar la evolución de su temperatura en el tiempo. Ahora bien, la información de SPAM no es una serie temporal, sino más bien una captura de áreas de cultivo en un momento determinado del tiempo. Con todo, estos datos son igualmente relevantes para capturar con precisión el estrés térmico sobre la producción del café.

Así, en una primera etapa del procesamiento de datos, se juntan ambas bases para obtener la evolución diaria de temperaturas mínimas y máximas dentro de cada grilla de 10 km x 10 km de un país específico. Luego, se toma el promedio simple de estas grillas para determinar una observación diaria por país, desde 1965 hasta 2019\footnote{Otra alternativa es ponderar por el área de cultivo --medido en hectáreas-- dentro de cada grilla. Sin embargo, dado que la base es equivalente a una foto en el tiempo, una ponderación estaría omitiendo posibles cambios en las zonas de cultivo en el tiempo, por lo que se prefiere utilizar un promedio simple.}.

La segunda etapa toma el insumo intermedio generado y calcula el tiempo de exposición que sufren los cultivos a temperaturas extremas, definidas por umbrales en grados centígrados. Así, en base a los antecedentes revisados en la sección anterior, se considera exposición a temperaturas bajas a aquellas bajo los 10°C; por su parte, se define como temperaturas en extremo altas a aquellas sobre los 30°C.

La literatura científica relativa al efecto de la temperatura sobre los cultivos de bienes agrícolas propone diferentes metodologías para la medición [CITAR]. En primer lugar, el uso de variables como \textit{Growing Degree Days} (GDD) posibilita medir la acumulación de grados centígrados dentro de cierto rango de temperaturas óptimas para los cultivos, bajo el supuesto de que la temperatura sigue una trayectoria diaria similar a una función sinusoidal [CITAR]. Otros trabajos adoptan aproximaciones más sencillas, en base a la diferencia entre la temperatura diaria promedio y un umbral de base sobre el cual el desarrollo de los cultivos alcanza condiciones óptimas [CITAR], encontrando leves diferencias cuando los datos contienen mucho ruido [CITAR].

Esta investigación propone otra alternativa, a fin de capturar esencialmente el tiempo de exposición a temperaturas extremas. En la Figura N°[] es posible comparar las metodologías revisadas. Mientras que, en Panel A, se captura la acumulación de calor en grados centígrados por día, el Panel B muestra la fracción del día en que los cultivos están expuestos a temperaturas sobre cierto umbral.

[FIGURA]

En un enfoque sencillo, la metodología propuesta considera que la temperatura sigue una distribución triangular cada día, comenzando y terminando con un valor mínimo, y presentando una temperatura máxima en el vértice superior. Aunque sencilla, esta aproximación permite obtener un valor directo para la fracción del día en que se alcanza un nivel de temperatura específico. Así, para un determinado umbral de temperatura $K$, y temperaturas mínimas y máximas ($t_{min}$ y $t_{max}$, respectivamente, con $t_{min} \leq t_{max}$), se tienen tres escenarios posibles, presentados en la Figura N°[]:

\begin{itemize}
    \item $t_{max} \leq K$: en este caso, la trayectoria diaria de la temperatura está siempre por debajo del umbral. Por lo tanto, si el objetivo es medir la fracción del día en que determinada zona geográfica estuvo expuesta a temperaturas en extremo bajas ($K = 10$), entonces la variable en construcción debe tomar un valor igual a 1 en este caso. Por otro lado, si el objetivo es capturar la exposición a temperaturas en extremo altas ($K = 30$), entonces la variable toma un valor igual a 0.
    
    \item $t_{min}  < K < t_{max}$: en este caso intermedio, la temperatura cruza el umbral en su trayectoria, lo que implica que la exposición a temperaturas en extremo altas o bajas solo cubre una fracción del día. Se puede demostrar que, si el objetivo es medir la exposición a temperaturas bajas, la variable toma un valor equivalente a $\frac{K - t_{min}}{t_{max} - t_{min}}$. En cambio, si lo que busca es cuantificar la exposición a temperaturas altas, entonces la variable toma el valor $1 - \frac{K - t_{min}}{t_{max} - t_{min}}$.
    
    \item $t_{min} \geq K$: este escenario es la contraparte del primero. La variable en cuestión puede tomar un valor igual a 1 o 0 según si se busca capturar la exposición a temperaturas altas ($K = 30$) o bajas ($K = 10$).
    
\end{itemize}

[FIGURA]

La siguiente tabla resume la información estadística de cada variable mencionada.

[TABLA]

\begin{table}[H]
\resizebox{\textwidth}{!}{%
\begin{tabular}{llccccc}
\hline
Variable                           & Unidad de medida                 & Observaciones & $\bar{x}$ & $\sigma$ & mín & máx \\ \hline
Precio ($P_{t}$)                   & Centavos de dólar por kg         &        658       &     224.79      &     105.61     & 77.07    & 694.37    \\
Exportaciones ($Q_{i,t}$)            & Millones de bolsas de 60kg       &      658         &     74.93      &    25.66      & 47.61    & 121.27    \\
PIB Importadores ($X_{i,t}$)         & Logaritmo de millones de dólares &      658         &     17.62      &     0.54     & 16.60    & 18.53    \\
Precio del té ($Z_{t}$)            & Centavos de dólares por kg       &        658       &     174.95      &    64.87      & 70.74    & 329.43    \\
Remuneraciones ($F_{i,t}$)           & Centavos de dólar por kg         &      658         &     144.87      &     88.18     & 0.78    & 677.58    \\
Choques climáticos ($\mathcal{W}$) & Variables 0/1                    &      658       & -         & -        & 0   & 1   \\ \hline
\end{tabular}}
\caption{Estadística descriptiva de las variables utilizadas en el análisis.}
\label{tab:4_1}
\end{table}